%%%%%%%%%%%%%%%%%%%%%%%%%%%%%%%%%%%%%%%%%%%%%%%%%%%%%%%%%%%%%%%%%%%%%%%%%%
% Research Scientist / Engineer -- early-stage AI company (ANONYMOUS
% recruiter posting, company not named). Founder with research background.
% Product: systems that get smarter about their users, learning from how
% someone interacts with the product and shaping the experience they get.
% Competitive salary + equity. Formative stage, real input on core tech.
% THE WORK (this is the spine of the tailoring):
%  1. Measure whether the system responds APPROPRIATELY given what it knows
%     about a given user  -> EVALUATION DESIGN
%  2. Turn REAL USAGE AND FEEDBACK into STRUCTURED SIGNALS the model can
%     learn from  -> candidate's "terminal impact" metric is exactly this
%  3. Run TRAINING EXPERIMENTS improving response quality based on HUMAN
%     FEEDBACK  -> partial: SFT yes, RLHF/RFT NOT claimed
%  4. Dig into FAILURE CASES: misreads a user, holds onto outdated info,
%     responds inconsistently  -> failure-mode & drift analysis
%  5. Weigh in on technical decisions as architecture comes together
% IDEAL BACKGROUND: a few years shipping ML/AI products; exposure to
% fine-tuning, evaluation, or human-feedback data; Python + deep learning
% tooling; self-starter comfortable without structure.
% CANDIDATE ANGLE: evaluation is his signature (evals, QoS, groundedness,
% drift, failure modes, 13/13 eval suite, 0.76 macro QoS); "terminal impact"
% is a genuine structured signal derived from real client conversations;
% early-stage hire = self-starter without structure; fine-tuning (LoRA,
% BERT) + PyTorch/HF; transformer + training loop from scratch = research
% depth for a research-led founder.
% HONEST -- NOT claimed: RLHF / RFT / reward modelling / preference data
% training, personalization or user-modelling systems, memory architectures
% (session/long-term/semantic), recommender systems, A/B testing,
% published research papers.
% Font: Source Sans Pro. Accent: teal + purple. Links: clickable labels.
% Compile: pdflatex (2 passes).
%%%%%%%%%%%%%%%%%%%%%%%%%%%%%%%%%%%%%%%%%%%%%%%%%%%%%%%%%%%%%%%%%%%%%%%%%%

\documentclass[10pt,a4paper]{article}

\usepackage[T1]{fontenc}
\usepackage[utf8]{inputenc}
\usepackage[default]{sourcesanspro}   % professional, neutral sans (Source Sans Pro)
\usepackage[a4paper,top=0.7cm,bottom=0.6cm,left=1.4cm,right=1.4cm]{geometry}
\usepackage{titlesec}
\usepackage{enumitem}
\usepackage{xcolor}
\usepackage{hyperref}
\usepackage{microtype}

% ---- Colours (echo the HTML CV: teal primary, purple accent) ----
\definecolor{accent}{HTML}{117A8B}   % teal  (hsl 187,74%,32%)
\definecolor{company}{HTML}{6C2BB3}  % purple (hsl 270,70%,45%)
\definecolor{muted}{HTML}{555555}

% ---- Links: clickable label text, coloured teal (no raw URLs) ----
\hypersetup{colorlinks=true, urlcolor=accent, linkcolor=accent}

% ---- Remove paragraph indent, set sensible line spacing ----
\setlength{\parindent}{0pt}
\linespread{0.99}
\pagestyle{empty}

% ---- Section headings: uppercase, bold, teal, thin rule underneath ----
\titleformat{\section}
  {\normalfont\large\bfseries\color{accent}\scshape}
  {}{0em}{}[{\color{accent!35}\titlerule[1pt]}]
\titlespacing*{\section}{0pt}{4pt}{3pt}

% ---- Compact bullet list ----
\newlist{tight}{itemize}{1}
\setlist[tight]{leftmargin=1.3em,itemsep=0.5pt,topsep=0.5pt,parsep=0pt,label=\textcolor{accent}{\small$\bullet$}}

% ---- Entry: company / dates on line 1, role / location on line 2 ----
\newcommand{\entry}[4]{%
  {\color{company}\bfseries #1}\hfill {\color{muted}#2}\par
  \vspace{1pt}
  {\itshape #3}\hfill {\color{muted}\itshape #4}\par
  \vspace{3pt}%
}

% ---- Inline divider for the contact line ----
\newcommand{\sep}{\,\textcolor{accent!60}{\textbar}\,}

\begin{document}

%======================= HEADER =======================
\begin{center}
  {\fontsize{24}{26}\selectfont\bfseries Abhijeet Lokhande}\\[3pt]
  {\color{accent}\large Research Engineer: Evaluation, Feedback Signals \& Applied ML}\\[5pt]
  {\small
    London, United Kingdom \sep
    \href{mailto:abhijeet_lokhande@proton.me}{abhijeet\_lokhande@proton.me} \sep
    07480801382 \\[2pt]
    \href{https://www.linkedin.com/in/ablds/}{LinkedIn} \sep
    \href{https://github.com/abhijeetscode}{GitHub} \sep
    \href{https://huggingface.co/ab-ai}{Hugging Face}
  }
\end{center}
\vspace{2pt}

%======================= SUMMARY =======================
\section{Summary}
Engineer who works on the hardest question in applied AI: whether a system is actually responding well, and how to turn that into something a model can learn from. As an early hire at a startup I have shipped LLM products to real users and built the measurement layer underneath them, including ``terminal impact,'' a metric I defined and computed from real client conversations that classifies the state a user is actually in when they leave a dialogue, because a completed conversation is not the same as a helped one. I spend much of my time on the failure cases: where a system misreads intent, drifts, or answers inconsistently, and what signal catches that class of error next time. I fine-tune models with PyTorch and the Hugging Face ecosystem, publish open models with 67K+ downloads, and I have built a GPT-style transformer and its training loop from first principles. Comfortable with little structure; that is how I have worked for four years.

%======================= SKILLS =======================
\section{Skills}
\begin{tight}
  \item[\textcolor{accent}{\small$\bullet$}] \textbf{Evaluation \& Measurement:} evaluation design \& eval suites, response-quality metrics, groundedness checks, quality-of-service scoring, failure-mode taxonomies, drift \& inconsistency detection, benchmark design, offline \& production evaluation
  \item[\textcolor{accent}{\small$\bullet$}] \textbf{Feedback \& Signals:} deriving structured signals from real usage and conversation data, outcome metrics from live user interactions, labelled evaluation sets, error analysis, instrumentation for feedback loops (OpenTelemetry, Arize Phoenix)
  \item[\textcolor{accent}{\small$\bullet$}] \textbf{Modelling \& Training:} PyTorch, TensorFlow, Hugging Face ecosystem, supervised fine-tuning (LoRA, BERT), embeddings \& representation learning, transformers built from scratch (training loop, checkpointing, GPU acceleration), GANs, Graph Convolution Networks
  \item[\textcolor{accent}{\small$\bullet$}] \textbf{Applied AI Products:} production LLM applications \& agentic systems, RAG \& semantic search, prompt \& context engineering, orchestration (LangGraph, LangChain), guardrails \& grounding gates
  \item[\textcolor{accent}{\small$\bullet$}] \textbf{Engineering:} Python (production-grade, TDD), Rust, SQL, FastAPI, PostgreSQL, Elasticsearch, Docker, CI/CD, AWS \& GCP, data pipelines over large messy real-world datasets
\end{tight}

%======================= EXPERIENCE =======================
\section{Experience}
\entry{Built AI (FinTech \& AI Company)}{April 2022 -- Present}{AI/ML Engineer (Early Stage Hire)}{London, UK}
\begin{tight}
  \item Built the measurement layer for production LLM systems serving 5+ leading enterprise customers: evaluation frameworks judging whether responses were appropriate and well grounded, quality-of-service metrics, and groundedness checks run before and after release.
  \item Turned real usage into structured signal: defined and computed ``terminal impact,'' an outcome metric derived from real client conversations that classifies whether a user exits a dialogue satisfied, unsatisfied, or unresolved, giving the team a learnable measure of whether conversations actually helped.
  \item Investigated failure cases as a discipline: where agents misread intent, relied on stale or ungrounded context, or answered inconsistently, using drift signals, failure-mode analysis, and incident-triggering alerts to catch whole classes of error.
  \item Fine-tuned and evaluated models with PyTorch and the Hugging Face ecosystem, and built the embedding and retrieval pipelines that shaped what a model saw before it answered.
  \item Shaped technical direction from the formative stage as an early hire, making architecture and tooling decisions with founders in an ambiguous, unstructured environment, and mentoring engineers who joined later.
  \item Shipped AI products end to end with measurable outcomes, including an agent that cut analyst modelling time by 70\% and an 85\% performance gain from rearchitecting a compute-heavy engine.
\end{tight}

\entry{King's College London}{March 2021 -- January 2022}{Research Associate}{London, UK}
\begin{tight}
  \item Ran deep-learning research experiments (Graph Convolution Networks, GANs) over a million-row dataset, achieving a 63\% improvement in the novelty score of generated structures, with a reproducible preprocessing and training pipeline.
\end{tight}

%======================= PROJECTS =======================
\section{Open-Source \& Projects}
\begin{tight}
  \item \textbf{Agentic Enterprise Assistant (evaluation-led):} LLM agent over grounded data with a grounding-validation gate checking every answer against tool results, role-based access, and auditable traces; OpenTelemetry to Arize Phoenix. Passed 13/13 evaluation cases across tool selection, grounding, and access compliance. \href{https://github.com/abhijeetscode/Agentic-Enterprise-Assistant}{View on GitHub}
  \item \textbf{Transformer from Scratch in Rust:} GPT-style decoder-only transformer built from first principles (multi-head causal self-attention, pre-norm blocks, full training loop with early stopping and checkpointing) on Candle with Metal GPU acceleration. \href{https://github.com/abhijeetscode/transformer-rust}{View on GitHub}
  \item \textbf{PII Detection LLMs:} fine-tuned GPT (LoRA) and BERT, published as open models; collectively 67K+ downloads on HuggingFace. \href{https://huggingface.co/ab-ai/pii_model}{View on Hugging Face}
\end{tight}

%======================= EDUCATION =======================
\section{Education}
{\color{company}\bfseries King's College London}\sep{\itshape MSc in Data Science}\hfill {\color{muted}London, UK}\par\vspace{1pt}
{\color{company}\bfseries C-DAC}\sep{\itshape PG Diploma in Advanced Computing (PG-DAC)}\hfill {\color{muted}Pune, India}\par\vspace{1pt}
{\color{company}\bfseries University of Pune}\sep{\itshape B.E. in Computer Science}\hfill {\color{muted}Pune, India}\par

\end{document}
